\section{A selective, borrow-tolerant ten-integer equation}
\label{sec:selective-sun}

There is a completely fixed polynomial in ten unrestricted integer unknowns
\[
 F_{10}(f,g,h,k,l,w,x,y,z,m)
\]
of exact total degree $21\,111\,400$, with a division-free evaluation circuit
using $659$ binary operations, for which ordinary mathematics formalizable
in ZFC proves
\[
 \neg\exists(f,g,h,k,l,w,x,y,z,m)\in\mathbb Z^{10}\;(F_{10}=0)
 \quad\Longleftrightarrow\quad\operatorname{Con}(\mathrm{ZFC}).
\]
The proof uses the pinned ProveIt machine/compiler calibration and Sun's
arithmetic bridge. It is not an end-to-end proof-assistant certificate and
it does not establish either consistency or inconsistency. The unknown
count is inherited, not a new universal unknown-count theorem.

Three source-layer changes are responsible for the improvement:
selective isolation of used quadratic monomials, one-bit target padding
that tolerates incoming borrows, and calibration by the actual mask
population instead of a uniform worst-case carry threshold. The latter
two changes are proved below and are independent of any search heuristic.
All scalar powers in this section have fixed nonnegative integer exponents.
Auxiliary letters are polynomial abbreviations, not extra unknowns.

\subsection{Fixed data and selective extraction}

Use precisely the 22 quadratic residuals of
\texttt{data/parent\_quartic31.json} and the complete fixed port tuple
\[
 (\beta_{\rm native},K_0,M_C,M_F)
 = (\texttt{beta},\texttt{K},\texttt{MC},\texttt{MF\_source}).
\]
The last port is already shifted; it must not be shifted again. Its tuple
is Wang943 $121\times8801$, with port bit lengths
$(1064921,1070004,1064921,1064922)$. Each payload is included, with compressed
and uncompressed hashes. The active binding is \texttt{INSTANCE.json};
historical relative paths inside preserved input files are provenance only.

Replace every positive parent coordinate by $n_i+1$. Let $r_j(n)$ denote
the resulting residual, written $r_j=\sum_\alpha p_{j,\alpha}n^\alpha$,
$|\alpha|\le2$. This is the original package-31 source, not the later
package-45 transport variant. The inherited source theorem is
\[
 \exists n\in\mathbb N^{31}\;\bigwedge_{j=1}^{22}r_j(n)=0
 \quad\Longleftrightarrow\quad\neg\operatorname{Con}(\mathrm{ZFC}),
 \qquad \mathbb N=\{0,1,2,\ldots\}.
\]
Its parent coordinate order is
\[
\begin{split}
 &(J,F,\alpha,z,f,h,i,o,j,s,w,\tau,\eta,\zeta,y,\gamma,\\
 &\hspace{25mm}N,E,P,c,a,X,D,C_2,L,F_2,V,g,I,B_{\rm aux},\Delta).
\end{split}
\]
Assign these coordinates the following weights, with a padding weight $a_0=0$:
\[
\begin{split}
(a_0,\ldots,a_{31})={}&(0,9,25,23,139,35,151,168,21,91,181,94,52,150,180,7,\\
 &135,38,88,183,106,178,136,20,138,8,64,105,154,37,137,120).
\end{split}
\]
The residual order and corresponding targets are
\[
\begin{split}
 &(16,2,9,14,21,7,15,22,13,3,19,20,1,10,17,8,4,6,11,18,12,5),\\
 &(70,86,451,685,1050,1296,1530,1874,2212,2358,2696,2780,\\
 &\hspace{7mm}2792,3157,3290,3634,3929,4238,4384,4527,4865,5190).
\end{split}
\]
Write $q_j$ for the target assigned to residual $r_j$, and put
\[
 t=183,\quad q=5190,\quad K_c=q+2t=5556,\quad E_c=t+K_c+2=5741,
 \quad H_c=2t+K_c+3+22=5947.
\]
For every used $\alpha$ and every $\gamma$ with $|\gamma|\le2$, the
finite certificate verifies
\begin{equation}\label{eq:sel-isolation}
 q_j-\alpha\cdot a+\gamma\cdot a=q_k
 \quad\Longleftrightarrow\quad j=k\ \hbox{and}\ \alpha=\gamma.
\end{equation}
All reversed exponents are nonnegative. There are 136 used source terms,
528 possible padded quadratic monomials, and $136\cdot528=71808$ checked
interactions. Only 267 different pair sums occur: a full Sidon condition
is deliberately not required.

Define
\[
 c_{j,\alpha}=(2-|\alpha|)!\prod_i\alpha_i!\;p_{j,\alpha},\qquad
 D_c(Z)=\sum_{j,\alpha}c_{j,\alpha}Z^{q_j-\alpha\cdot a}.
\]
Condition \eqref{eq:sel-isolation} implies that the 136 displayed exponents
in $D_c$ are distinct. Indeed, if two were equal, multiplying by the
monomial belonging to either source term would violate isolation at its
target. Thus, for $A_c\ge\max|c_{j,\alpha}|$, every coefficient of $D_c$
has absolute value at most $A_c$, with no multiplicity factor.
Its leading coefficient is exactly $2$: the last residual is $r_5$ and
$r_5(0)=1$. In particular, the zero vector is not a source solution.

For $C_n(Z)=(1+\sum_i n_iZ^{a_i})^2$, multinomial expansion gives
\begin{equation}\label{eq:sel-extraction}
 [Z^{q_j}]C_n(Z)D_c(Z)=2r_j(n),\qquad
 \left|[Z^s]C_nD_c\right|\le A_c(1+\textstyle\sum_i n_i)^2.
\end{equation}
The factorial multiplier cancels the multinomial coefficient; isolation
excludes every unwanted contribution. This proves the identity for all
source vectors, not just tested examples.

Take $A_c$ to be the maximum coefficient majorant obtained by replacing
in each fixed-port recipe a monomial by $2^{\sum e_i\ell_i}$, multiplying
by the absolute value of its integer coefficient, and summing. Here
$\ell_i$ are the four verified port bit lengths. Exact integer arithmetic
gives $\operatorname{bitlength}(A_c)=2134927$ and, with
\[
 \kappa=2134940,\qquad \beta_S=2^\kappa,
 \qquad \left\lfloor\beta_S/A_c\right\rfloor=16383,
\]
we have
\begin{equation}\label{eq:sel-bound}
 \beta_S>4096A_c\quad\hbox{and}\quad\beta_S>64A_c.
\end{equation}

\subsection{The new source polynomials}

In the two unrestricted integer unknowns $f,g$, define
\begin{align*}
 b&=1+3f,& B_c&=\beta_S b^2,& \mathcal C&=bB_c^t,\\
 \ell&=\sum_{i=1}^{31}B_c^{a_i},&
 N_0&=B_c^{t+1},&N_1&=B_c^{K_c+1},&N&=N_0N_1,\\
 A_T&=\sum_{j=1}^{22}B_c^{q_j},&
 M_0&=N_0-1-(b-1)\ell,\\
 T_c&=M_0+(B_c-2)A_TN_0,&
 J_c&=2(1+g)^2D_c(B_c)+A_T,& S_c&=g+J_cN_0,\\
 X&=N^2+(S_c+T_c)N+N-1-S_c,\\
 Y&=\beta_S^{15}\,2^{\kappa-21}\,B_c^{5931}b.
\end{align*}
Equivalently, $Y=2^{\kappa H_c-21}b^{11863}$; the first form avoids
materializing its larger fixed coefficient. The fixed leaf
$2^{\kappa-21}=2^{2134919}$ is explicitly specified, not an unknown.
Also $\mathcal C=2^{\kappa t}b^{367}$ has the monomial form required
by the inherited bridge.

\paragraph{Bounds before binary-base recovery.}
Suppose $f\ne0$, $b$ is an integer square and $0\le g<2\mathcal C$.
Since $b\equiv1\pmod3$, these hypotheses imply $b\ge4$ and $f>0$.
The integer $B_c$ is positive and even, but is not yet assumed to be a
power of two. Its size implies $B_c>2b$ and $B_c>A_c+1$.
In the following estimates write $B=B_c$ temporarily.

The base-$B$ digits of $M_0$ are $B-1$ except at the 31 designated
positions, where they are $B-b$. Hence $0\le M_0<N_0$.
The leading term of $D_c(B)$ is $2B^q$, and
\[
 |D_c(B)-2B^q|\le A_c\frac{B^q-1}{B-1}<B^q,
 \qquad 0<D_c(B)<\frac{A_c B^{q+1}}{B-1}.
\]
Because $g$ is integral, $1+g\le2bB^t$. Therefore
\[
 0<2(1+g)^2D_c(B)
 <\frac{8A_cb^2}{B-1}N_1<\frac14N_1.
\]
The last inequality follows from $\beta_S>64A_c$ and $B>2$.
Also $(B-1)A_T<B^{q+1}\le N_1/B^2$, since $t\ge1$.
Consequently
\[
 0<J_c+(B-2)A_T<\tfrac12N_1.
\]
We have $g<N_0$, so $g+M_0<2N_0$. As $N_1\ge4$, this proves
\begin{equation}\label{eq:sel-tight-range}
 0\le S_c,\quad0\le T_c,\qquad S_c+T_c<N.
\end{equation}
In particular $X\ge N^2>3b$ and $Y>\max(b,256)$.
Furthermore $B_c\equiv N_0\equiv N\equiv1\pmod3$,
$M_0\equiv0$ and $T_c\equiv-22$, giving
\[
 X\equiv1+T_c\equiv0\pmod3.
\]
These statements were obtained without assuming that a square $b$ is a
power of two. This order matters in the reverse bridge.

\subsection{Sparse padding that tolerates borrows}

Now assume additionally that $b$ is a power of two. Then $B_c$ and $N_0$
are powers of two. The condition $g\mathbin{\&}M_0=0$ and $0\le g<N_0$
is exactly
\[
 g=\sum_{i=1}^{31}n_iB_c^{a_i},\qquad 0\le n_i<b.
\]
The weights are distinct, so the tuple is unique. Indeed, the low mask
forbids every bit outside the low $\log_2 b$ bits of these designated
base-$B_c$ cells.

For this tuple let $h_s=[Z^s]C_nD_c$ and let $\epsilon_s$ be the indicator
of the 22 target positions. Equations \eqref{eq:sel-extraction} and
\eqref{eq:sel-bound} give $|h_s|<B_c/4$, and
\[
 J_c=\sum_{s=0}^{K_c}(2h_s+\epsilon_s)B_c^s.
\]
This is a signed-coefficient expansion, not yet a digit expansion.
Set $u_0=0$ and carry out ordinary Euclidean digit normalization:
\[
 u_{s+1}=\left\lfloor\frac{2h_s+\epsilon_s+u_s}{B_c}\right\rfloor,
 \qquad d_s=2h_s+\epsilon_s+u_s-B_cu_{s+1}.
\]
Induction gives $u_s\in\{-1,0\}$ and $0\le d_s<B_c$.
At a marked position $\epsilon_s=1$, the binary mask $B_c-2$ permits
exactly the digits $0$ and $1$. If $h_s=0$, the unnormalized digit is
$1+u_s\in\{0,1\}$, so it passes. If $h_s\ge1$, that digit lies in
$[2,B_c)$ and fails. If $h_s\le-1$, it lies strictly between $-B_c+1$
and $0$, so its normalized digit is at least $2$ and again fails.
Thus each marked test detects precisely $h_s=0$, regardless of an
incoming borrow. Positivity of $J_c$ was already proved, so these are
ordinary digits of a nonnegative integer.

Concatenation at the power-of-two boundary $N_0$ now yields
\begin{equation}\label{eq:sel-carry-source}
 S_c\mathbin{\&}T_c=0
 \quad\Longleftrightarrow\quad
 g=\sum_i n_iB_c^{a_i},\ 0\le n_i<b,\quad r_1(n)=\cdots=r_{22}(n)=0.
\end{equation}
The dense geometric padding used by the previous construction has been
replaced by only one extra bit at each marked cell. Unmarked negative
coefficients may generate borrows; the proof allows them rather than
silently assuming them absent.

\subsection{Calibrating the binomial threshold by the mask population}

For $N=2^u$ and nonnegative integers $S,T$ satisfying $S+T<N$, set
\[
 X_*=N^2+(S+T)N+N-1-S.
\]
Its leading, middle, and low blocks are disjoint. If $s_2$ denotes binary
population, then
\begin{align*}
 s_2(X_*)&=1+s_2(S+T)+u-s_2(S)\\
 &=1+u+s_2(T)-\bigl(s_2(S)+s_2(T)-s_2(S+T)\bigr).
\end{align*}
The parenthesized term is nonnegative and vanishes exactly when
$S\mathbin{\&}T=0$. Legendre's formula gives
$\nu_2\binom{2X_*}{X_*}=s_2(X_*)$. Consequently
\begin{equation}\label{eq:sel-popcount}
 2^{1+u+s_2(T)}\mid\binom{2X_*}{X_*}
 \quad\Longleftrightarrow\quad S\mathbin{\&}T=0.
\end{equation}
The stronger range $S+T<N$ is essential and was proved in
\eqref{eq:sel-tight-range}; separate inequalities $S<N,T<N$ would not suffice.

Write $b=2^k$ and $B_c=2^v$ on the binary branch. The low mask has
population $(t+1)v-31k$, while each of the 22 high marked cells contributes
$v-1$ bits. Hence
\[
 s_2(T_c)=(t+1+22)v-31k-22,\qquad u=E_cv,
\]
and
\[
 1+u+s_2(T_c)=H_cv-31k-21
 =\kappa H_c+(2H_c-31)k-21=\log_2Y.
\]
Combining \eqref{eq:sel-popcount} and \eqref{eq:sel-carry-source} proves
the required reverse source implication. Conversely, a source solution
is nonzero because $r_5(0)=1$. Choose arbitrarily large $b=4^s$ exceeding
all its coordinates and put $f=(b-1)/3$. Its encoding satisfies
$b\le B_c\le g<\mathcal C$ and all carry tests pass. Thus it supplies
arbitrarily large forward witnesses with $Y\mid\binom{2X}{X}$.
This establishes the entire source contract for the unchanged finalizer.

\subsection{Explicit radical-free finalizer}

The following formulas are the inherited package-31 quintic,
merged-square, structural-weight finalizer. They are written out to fix
the actual equation. Distinguish its local $B_*,S_*,T_*$ from the source
base and masks above. Define
\begin{gather*}
 L=lY,\qquad U=2LX,\qquad V=4gwY,\qquad W=bw,\\
 A=U(V+1),\qquad B_*=2X+1,\qquad C=B_*+(A-2)h,\\
 P=U^2V,\qquad K=X+1+k(P-2),\\
 S_*=2A-5,\qquad T_*=3WC-2(W^2-1),\\
 R_2=(P^2-4)K^2+4,\\
 \Delta=8\mathcal C^3gK^2-g^2\{32(C-KL)^2\mathcal C^3+g^2K^2\},\\
 O=f^2l^2\Delta,\qquad \xi=zx,\\
 D=(A^2-4)C^2+4,\qquad E=C^2D\xi,\\
 F=4(A^2-4)E^2+1,\qquad G_*=2(A+2)(A-2)^2E^2-1,\\
 \Omega=16E,\qquad H_*=C+(B_*-C)F+\Omega Fy,\\
 I_*=(G_*^2-1)H_*^2+1.
\end{gather*}
Put
\begin{align*}
 r&=T_*+S_*y,&\Sigma&=\xi S_*^5,\\
 \Theta&=(r^3-rS_*^2-S_*^3)(r^2+23S_*^2),\\
 a_0&=b\{1+(3O-4)b^2z^2\},&a_1&=R_2,&a_2&=DFI_*,\\
 w_0&=1,&w_1&=a_0,&w_2&=a_0(P^2+K^2+3),\\
 Z&=\Sigma m-\Theta,&A_i&=\Sigma^2w_i^2a_i\quad(0\le i\le2),\\
 V_n&=Z^2+A_2-A_0-A_1,&W_n&=V_n^2+4Z^2A_2-4A_0A_1.
\end{align*}
The promised single equation is
\begin{equation}\label{eq:sel-final-ten}
 \boxed{\ F_{10}=W_n^2-16V_n^2Z^2A_2=0.\ }
\end{equation}
Every expression is an integer polynomial. In particular, radicals and
binomial coefficients appear in the proof only, never in the equation.

\paragraph{Exact dependency boundary.}
Sun's 2021 Theorems 4.12--4.13, specialized as in his 2026 Propositions
6.1--6.2, provide the two bridge directions for any source satisfying the
contract just proved. Forward witnesses have $b$ square, $R_2$ square,
$S_*\mid T_*$, $C=u_{B_*}(A)$ and $O>0$. Reverse witnesses with
$b$ square, $O>0$, $\xi\ne0$, $R_2,DFI_*$ square and $S_*\mid T_*$
recover $b\in2\uparrow$, $0<g<2\mathcal C$ and the required binomial
divisibility. The residue interpolation theorem permits $\xi$ to be
divisible by any prescribed positive integer and permits any residue
class for $y\pmod\xi$.
The only data substituted here are the newly proved source polynomials;
the bridge does not require $Y=N^2$.

For completeness, the inherited finalizer implications can be verified
as follows. The polynomial $(v^3-v-1)(v^2+23)$ has a root modulo every
positive integer and none over $\mathbb Q$. For odd primes other than
23 the cubic discriminant is $-23$: if the cubic is irreducible its
Frobenius permutation is even, so the quadratic has a simple root;
otherwise the cubic has a simple root. At 23 the cubic has simple root
3. At 2, $-23\equiv1\pmod8$. Hensel lifting and the Chinese remainder
theorem finish the local-root argument. Prime valuations in the two
homogeneous factors show that $\Sigma\mid\Theta$ forces $S_*\mid T_*$:
if $v_p(S_*)>v_p(T_*)$, the valuation of $\Theta$ is
$5v_p(T_*)<5v_p(S_*)$. Since $S_*$ is odd, it never vanishes; the absence
of rational roots ensures $\Theta\ne0$.

When $z\ne0$, $a_0$ square implies that its coprime factors $b$ and
$1+(3O-4)b^2z^2$ are signed squares. The congruence $b\equiv1\pmod3$
excludes a negative square and excludes zero, including the zero-factor
exception. Thus $b$ is a positive square and the usual positivity Pell
lemma gives $O>0$. Conversely, forward witnesses have $O\ge4$: the
positive integer $\Delta$ is divisible by $g\ge b\ge4$. For
$d=3O-4>0$, $d\equiv2\pmod3$ is nonsquare. A positive Pell unit raised
to its order modulo $b$ gives $v=bz>0$ with $u^2-dv^2=1$, making $a_0$
a square. Residue interpolation then supplies $\xi=zx$ and a $y$ that
simultaneously makes $a_2$ square and $\Sigma\mid\Theta$.

If $z\ne0$, the displayed $a_0$ always has $|a_0|>1$. The exceptional
possibilities $a_0=0,\pm1$ force $|b|=1$, hence $b=1$, $f=O=0$, and
$a_0=1-4z^2$, a contradiction. Also
$\sqrt{|a_1|}\le P^2+K^2+2$. Thus the weights obey
\[
 |w_1|>\sqrt{|a_0|},\qquad
 |w_2|>\sqrt{|a_0|}+|w_1|\sqrt{|a_1|}.
\]
In the multiquadratic field, an automorphism changing a nonzero weighted
root has a largest changed index; its magnitude exceeds the sum of all
smaller changed magnitudes. A rational weighted sum can therefore occur
only when every $a_i$ is an integer square. This also covers negative
radicands by complex absolute values, and zero radicands separately.

The exact norm identity
\[
 F_{10}=\prod_{\epsilon\in\{\pm1\}^3}
 \left(Z+\sum_{i=0}^2\epsilon_i\sqrt{A_i}\right)
\]
has 35 expanded terms and is checked independently by integer polynomial
arithmetic. If $\Sigma=0$ then $F_{10}=\Theta^8\ne0$. At a zero of
$F_{10}$ we therefore have $\Sigma\ne0$, hence $z\ne0$; separation forces
all $a_i$ square and $\Sigma\mid\Theta$. The reverse bridge and source
proof recover a parent solution. Conversely the forward construction
just described makes all $a_i$ square and $\Sigma\mid\Theta$; choosing
$m$ to cancel one signed factor gives $F_{10}=0$. Composition with the
pinned parent theorem proves the stated consistency equivalence in ZFC.

\subsection{Exact finite certificates and comparison}

The literal DAG \texttt{generated/signed10\_b2.json} has ten variable
leaves, 18 fixed leaves, and 659 live binary arithmetic gates:
214 additions, 26 subtractions and 419 multiplications. Multiplication
by a numeral is charged; fixed-coefficient construction is outside this
arithmetic-circuit cost. It is not a bit-complexity claim.

The independent checker does not import or execute the builder. It
reconstructs the shifted quadratic source and fixed recipes, verifies
all selective interactions, checks 52 local polynomial circuit identities
exactly, checks the norm identity, and audits liveness and the gate ledger.
For the total degree it propagates a syntactic upper bound, then tracks
the coefficient of that degree after substituting the same indeterminate
for all ten variables. The resulting nonzero residues are
\[
 \begin{array}{c|ccc}
 p&101&1009&1000003\\\hline
 [T^{21111400}]F_{10}(T,\ldots,T)\pmod p&36&42&374910.
 \end{array}
\]
Any one of these gives the exact lower bound, hence the total degree is
exactly $21\,111\,400$. This is not a random evaluation argument.
Finite small-base audits additionally corroborate the borrow and
population lemmas, including non-power-of-two square source bases.
They are tests of implementations, not substitutes for the unbounded proofs.

The previous package-40 ten-integer point was degree $237\,520\,456$
with 782 operations. The new degree is lower by a factor exceeding
$11.25$, and the circuit uses 123 fewer operations, a reduction of about
$15.7\%$. Package 43's 809-operation variant is also improved. These are
comparisons of explicit constructions at the pinned repository snapshot,
not claims of global optimality, of an absolute ``simplest'' equation, or
of improvement over different variable/domain tradeoffs.

The pinned revision is
\texttt{124c07c0b77a9e20e84676514d8b1ff78536ea3f}.
\texttt{SOURCE\_PINS.json} records the verified exact source bytes.
Primary arithmetic references are Z.-W. Sun,
\emph{Further results on Hilbert's Tenth Problem}, Sci. China Math. 64
(2021), 281--306, \url{https://arxiv.org/html/1704.03504v7}; and
\emph{Ten unknowns for Hilbert's tenth problem over the integers},
2026 preprint, \url{https://arxiv.org/html/2609.01594v2}.
The inherited finalizer proof and coefficient adapter are preserved in
\texttt{upstream/}. Their original relative links describe upstream
layouts; the source pin manifest provides their current verified origins.
